\documentclass[aspectratio=169, 11pt]{beamer}
% ============================================================================
% Preambulo compartido para presentaciones Beamer
% Estadistica 2026-II - Pontificia Universidad Javeriana
% Incluir con: % ============================================================================
% Preambulo compartido para presentaciones Beamer
% Estadistica 2026-II - Pontificia Universidad Javeriana
% Incluir con: % ============================================================================
% Preambulo compartido para presentaciones Beamer
% Estadistica 2026-II - Pontificia Universidad Javeriana
% Incluir con: \input{../preambulo_beamer.tex} despues de \documentclass
% ============================================================================

\usepackage[T1]{fontenc}
\usepackage[utf8]{inputenc}
\usepackage[spanish]{babel}
\usepackage{amsmath, amssymb, amsfonts}
\usepackage{booktabs}
\usepackage{graphicx}
\usepackage{tikz}
\usetikzlibrary{shapes.geometric, positioning, arrows.meta, calc}
\usepackage{pgfplots}
\pgfplotsset{compat=1.18}
\usepackage{listings}
\usepackage{xcolor}
\usepackage{hyperref}
\usepackage{multicol}


% --- Tema Beamer (minimalista) ---
\usetheme{default}
\usecolortheme{default}
\usefonttheme{default}
\setbeamertemplate{navigation symbols}{}
\setbeamertemplate{caption}[numbered]
\setbeamertemplate{footline}{%
  \hfill\usebeamerfont{footline}\insertframenumber/\inserttotalframenumber\hspace{1em}\vspace{0.5em}%
}
\setbeamertemplate{itemize items}[circle]
\setbeamertemplate{enumerate items}[default]

% --- Colores institucionales (sutiles) ---
\definecolor{javeriana}{RGB}{0, 62, 105}
\definecolor{javerianalight}{RGB}{0, 105, 170}
\definecolor{codigobg}{RGB}{245, 245, 245}
\definecolor{codigofg}{RGB}{40, 40, 40}
\definecolor{comentario}{RGB}{100, 100, 100}
\definecolor{keyword}{RGB}{0, 100, 150}
\definecolor{string}{RGB}{180, 60, 30}

\setbeamercolor{title}{fg=javeriana}
\setbeamercolor{frametitle}{fg=javeriana}
\setbeamercolor{structure}{fg=javeriana}
\setbeamercolor{block title}{fg=white, bg=javeriana!75}
\setbeamercolor{block body}{bg=javeriana!5}
\setbeamercolor{block title alerted}{fg=white, bg=red!70!black}
\setbeamercolor{block body alerted}{bg=red!5}
\setbeamercolor{block title example}{fg=white, bg=green!40!black}
\setbeamercolor{block body example}{bg=green!5}
\setbeamercolor{item}{fg=javeriana}

\setbeamerfont{title}{series=\bfseries, size=\Large}
\setbeamerfont{frametitle}{series=\bfseries}

% --- Configuracion de listings para R ---
\lstset{
  language=R,
  basicstyle=\ttfamily\footnotesize,
  backgroundcolor=\color{codigobg},
  commentstyle=\color{comentario}\itshape,
  keywordstyle=\color{keyword}\bfseries,
  stringstyle=\color{string},
  numbers=none,
  frame=single,
  rulecolor=\color{javeriana!30},
  breaklines=true,
  showstringspaces=false,
  columns=flexible,
  morekeywords={library, ggplot, aes, geom_histogram, geom_boxplot,
    geom_point, geom_smooth, geom_density, geom_bar, geom_line,
    geom_vline, geom_hline, geom_errorbar, labs, theme_minimal,
    facet_wrap, scale_fill_manual, scale_color_manual,
    summarise, group_by, filter, mutate, select, arrange,
    read.csv, head, str, summary, mean, median, sd, var, cor,
    lm, t.test, chisq.test, wilcox.test, kruskal.test, aov,
    pnorm, qnorm, dnorm, rnorm, qt, pt, qchisq, pchisq, qf, pf,
    confint, coeftest, vcovHC, bptest, vif, cooks.distance,
    TukeyHSD, table, prop.table}
}

% --- Comandos personalizados ---
\newcommand{\R}{\texttt{R}}
\newcommand{\code}[1]{\texttt{\footnotesize #1}}
\newcommand{\variable}[1]{\texttt{#1}}
\newcommand{\paquete}[1]{\texttt{#1}}

% --- Informacion del curso ---
\institute{Pontificia Universidad Javeriana\\
  Facultad de Ciencias Económicas y Administrativas}
\date{2026-II}
 despues de \documentclass
% ============================================================================

\usepackage[T1]{fontenc}
\usepackage[utf8]{inputenc}
\usepackage[spanish]{babel}
\usepackage{amsmath, amssymb, amsfonts}
\usepackage{booktabs}
\usepackage{graphicx}
\usepackage{tikz}
\usetikzlibrary{shapes.geometric, positioning, arrows.meta, calc}
\usepackage{pgfplots}
\pgfplotsset{compat=1.18}
\usepackage{listings}
\usepackage{xcolor}
\usepackage{hyperref}
\usepackage{multicol}


% --- Tema Beamer (minimalista) ---
\usetheme{default}
\usecolortheme{default}
\usefonttheme{default}
\setbeamertemplate{navigation symbols}{}
\setbeamertemplate{caption}[numbered]
\setbeamertemplate{footline}{%
  \hfill\usebeamerfont{footline}\insertframenumber/\inserttotalframenumber\hspace{1em}\vspace{0.5em}%
}
\setbeamertemplate{itemize items}[circle]
\setbeamertemplate{enumerate items}[default]

% --- Colores institucionales (sutiles) ---
\definecolor{javeriana}{RGB}{0, 62, 105}
\definecolor{javerianalight}{RGB}{0, 105, 170}
\definecolor{codigobg}{RGB}{245, 245, 245}
\definecolor{codigofg}{RGB}{40, 40, 40}
\definecolor{comentario}{RGB}{100, 100, 100}
\definecolor{keyword}{RGB}{0, 100, 150}
\definecolor{string}{RGB}{180, 60, 30}

\setbeamercolor{title}{fg=javeriana}
\setbeamercolor{frametitle}{fg=javeriana}
\setbeamercolor{structure}{fg=javeriana}
\setbeamercolor{block title}{fg=white, bg=javeriana!75}
\setbeamercolor{block body}{bg=javeriana!5}
\setbeamercolor{block title alerted}{fg=white, bg=red!70!black}
\setbeamercolor{block body alerted}{bg=red!5}
\setbeamercolor{block title example}{fg=white, bg=green!40!black}
\setbeamercolor{block body example}{bg=green!5}
\setbeamercolor{item}{fg=javeriana}

\setbeamerfont{title}{series=\bfseries, size=\Large}
\setbeamerfont{frametitle}{series=\bfseries}

% --- Configuracion de listings para R ---
\lstset{
  language=R,
  basicstyle=\ttfamily\footnotesize,
  backgroundcolor=\color{codigobg},
  commentstyle=\color{comentario}\itshape,
  keywordstyle=\color{keyword}\bfseries,
  stringstyle=\color{string},
  numbers=none,
  frame=single,
  rulecolor=\color{javeriana!30},
  breaklines=true,
  showstringspaces=false,
  columns=flexible,
  morekeywords={library, ggplot, aes, geom_histogram, geom_boxplot,
    geom_point, geom_smooth, geom_density, geom_bar, geom_line,
    geom_vline, geom_hline, geom_errorbar, labs, theme_minimal,
    facet_wrap, scale_fill_manual, scale_color_manual,
    summarise, group_by, filter, mutate, select, arrange,
    read.csv, head, str, summary, mean, median, sd, var, cor,
    lm, t.test, chisq.test, wilcox.test, kruskal.test, aov,
    pnorm, qnorm, dnorm, rnorm, qt, pt, qchisq, pchisq, qf, pf,
    confint, coeftest, vcovHC, bptest, vif, cooks.distance,
    TukeyHSD, table, prop.table}
}

% --- Comandos personalizados ---
\newcommand{\R}{\texttt{R}}
\newcommand{\code}[1]{\texttt{\footnotesize #1}}
\newcommand{\variable}[1]{\texttt{#1}}
\newcommand{\paquete}[1]{\texttt{#1}}

% --- Informacion del curso ---
\institute{Pontificia Universidad Javeriana\\
  Facultad de Ciencias Económicas y Administrativas}
\date{2026-II}
 despues de \documentclass
% ============================================================================

\usepackage[T1]{fontenc}
\usepackage[utf8]{inputenc}
\usepackage[spanish]{babel}
\usepackage{amsmath, amssymb, amsfonts}
\usepackage{booktabs}
\usepackage{graphicx}
\usepackage{tikz}
\usetikzlibrary{shapes.geometric, positioning, arrows.meta, calc}
\usepackage{pgfplots}
\pgfplotsset{compat=1.18}
\usepackage{listings}
\usepackage{xcolor}
\usepackage{hyperref}
\usepackage{multicol}


% --- Tema Beamer (minimalista) ---
\usetheme{default}
\usecolortheme{default}
\usefonttheme{default}
\setbeamertemplate{navigation symbols}{}
\setbeamertemplate{caption}[numbered]
\setbeamertemplate{footline}{%
  \hfill\usebeamerfont{footline}\insertframenumber/\inserttotalframenumber\hspace{1em}\vspace{0.5em}%
}
\setbeamertemplate{itemize items}[circle]
\setbeamertemplate{enumerate items}[default]

% --- Colores institucionales (sutiles) ---
\definecolor{javeriana}{RGB}{0, 62, 105}
\definecolor{javerianalight}{RGB}{0, 105, 170}
\definecolor{codigobg}{RGB}{245, 245, 245}
\definecolor{codigofg}{RGB}{40, 40, 40}
\definecolor{comentario}{RGB}{100, 100, 100}
\definecolor{keyword}{RGB}{0, 100, 150}
\definecolor{string}{RGB}{180, 60, 30}

\setbeamercolor{title}{fg=javeriana}
\setbeamercolor{frametitle}{fg=javeriana}
\setbeamercolor{structure}{fg=javeriana}
\setbeamercolor{block title}{fg=white, bg=javeriana!75}
\setbeamercolor{block body}{bg=javeriana!5}
\setbeamercolor{block title alerted}{fg=white, bg=red!70!black}
\setbeamercolor{block body alerted}{bg=red!5}
\setbeamercolor{block title example}{fg=white, bg=green!40!black}
\setbeamercolor{block body example}{bg=green!5}
\setbeamercolor{item}{fg=javeriana}

\setbeamerfont{title}{series=\bfseries, size=\Large}
\setbeamerfont{frametitle}{series=\bfseries}

% --- Configuracion de listings para R ---
\lstset{
  language=R,
  basicstyle=\ttfamily\footnotesize,
  backgroundcolor=\color{codigobg},
  commentstyle=\color{comentario}\itshape,
  keywordstyle=\color{keyword}\bfseries,
  stringstyle=\color{string},
  numbers=none,
  frame=single,
  rulecolor=\color{javeriana!30},
  breaklines=true,
  showstringspaces=false,
  columns=flexible,
  morekeywords={library, ggplot, aes, geom_histogram, geom_boxplot,
    geom_point, geom_smooth, geom_density, geom_bar, geom_line,
    geom_vline, geom_hline, geom_errorbar, labs, theme_minimal,
    facet_wrap, scale_fill_manual, scale_color_manual,
    summarise, group_by, filter, mutate, select, arrange,
    read.csv, head, str, summary, mean, median, sd, var, cor,
    lm, t.test, chisq.test, wilcox.test, kruskal.test, aov,
    pnorm, qnorm, dnorm, rnorm, qt, pt, qchisq, pchisq, qf, pf,
    confint, coeftest, vcovHC, bptest, vif, cooks.distance,
    TukeyHSD, table, prop.table}
}

% --- Comandos personalizados ---
\newcommand{\R}{\texttt{R}}
\newcommand{\code}[1]{\texttt{\footnotesize #1}}
\newcommand{\variable}[1]{\texttt{#1}}
\newcommand{\paquete}[1]{\texttt{#1}}

% --- Informacion del curso ---
\institute{Pontificia Universidad Javeriana\\
  Facultad de Ciencias Económicas y Administrativas}
\date{2026-II}


\title{Sesión 10: Regresión Lineal Múltiple --- OLS}
\subtitle{Anatomía de una investigación económica con datos experimentales}
\author{Prof.\ Jaime Polanco Jimenez}
\date{2026-II}

\begin{document}

\begin{frame}
\titlepage
\end{frame}

\begin{frame}{Agenda}
\small
\tableofcontents
\end{frame}

% ============================================================================
\section{Motivación: Una investigación real}
% ============================================================================

\begin{frame}{¿Por qué necesitan regresión múltiple?}
\begin{exampleblock}{Porque un solo predictor no basta}
\small
Tenemos datos de un experimento real (AI-Driven Mentorship, KU Leuven, 2024):
\begin{itemize}\setlength{\itemsep}{2pt}
  \item 148 estudiantes, 4 grupos: Control, Friendly, Mentor, Narrative
  \item Queremos saber si el tipo de mentor IA afecta la intención de inscripción universitaria
  \item Pero esa intención también depende de la intención \textit{previa}, autoeficacia, origen migrante\ldots
  \item Si no controlamos por estos factores, nuestras conclusiones serán erróneas
\end{itemize}
\end{exampleblock}

\begin{block}{La frase clave del curso}
\small
``Manteniendo constantes las demás variables'' (ceteris paribus). Esto es lo que distingue un análisis serio de uno superficial.
\end{block}
\end{frame}

\begin{frame}{La pregunta de investigación}
\begin{block}{Contexto}
\small
En Flandes (Bélgica), estudiantes de origen migrante tienen menor probabilidad de acceder a educación superior, \textbf{incluso controlando por desempeño académico previo}. ¿Por qué? Fallas de información: no conocen los programas, sobreestiman costos, subestiman retornos.
\end{block}

\vspace{0.2cm}
\begin{exampleblock}{Pregunta}
\small
¿Puede un chatbot con IA corregir estas fallas de información? Y más importante: ¿importa \textbf{cómo} se entrega la información? ¿Un tono ``cálido'' es mejor que uno puramente factual?
\end{exampleblock}

\vspace{0.2cm}
\textbf{La tensión central:}
\begin{itemize}
  \item Mecanismo \textbf{informacional}: la IA provee datos sobre costos, becas, programas
  \item Mecanismo \textbf{relacional}: la IA simula empatía, motivación, narrativas inspiradoras
  \item ¿Cuál domina? ¿Depende del tipo de estudiante?
\end{itemize}
\end{frame}

\begin{frame}{Asimetría de información y capital humano}
\begin{block}{Modelo de Capital Humano (HCM)}
Un estudiante se inscribe en educación superior si los beneficios esperados superan los costos:
\[
\sum_{t=1}^{T} \frac{\mathbb{E}[B_t \mid \mathcal{I}_i]}{(1+r)^t} \;>\; \sum_{t=0}^{T} \frac{\mathbb{E}[C_t \mid \mathcal{I}_i]}{(1+r)^t}
\]
donde $\mathcal{I}_i$ es la \textbf{información disponible} y $r$ es la tasa de descuento (paciencia).
\end{block}

\vspace{0.3cm}
\textbf{Pregunta clave:} ¿Qué pasa cuando $\mathcal{I}_i$ es incompleta?

Si las familias no tienen toda la información, sus decisiones educativas pueden ser \textbf{subóptimas} --- incluso cuando ir a la universidad les convendría.
\end{frame}

\begin{frame}{La cuña informacional}
\begin{alertblock}{Cuando $\mathcal{I}_i \subset \mathcal{I}^*$ (información incompleta)}
\small
Las familias migrantes enfrentan 4 fricciones informacionales:
\begin{enumerate}\setlength{\itemsep}{3pt}
  \item \textbf{Subestiman retornos} ($\mathbb{E}[B_t \mid \mathcal{I}_i] < B_t$): no observan graduados exitosos
  \item \textbf{Sobreestiman costos} ($\mathbb{E}[C_t \mid \mathcal{I}_i] > C_t$): desconocen becas y ayudas
  \item \textbf{Desajuste percibido}: no entienden el sistema de tracks (ASO, TSO, BSO)
  \item \textbf{Baja autoeficacia}: creen que no pueden tener éxito académico
\end{enumerate}
\end{alertblock}

\vspace{0.3cm}
\textbf{Resultado:} \textbf{no-inscripción subóptima} incluso cuando el retorno neto es positivo.

\vspace{0.2cm}
{\small Nuestro experimento ataca estas fricciones con 4 chatbots de IA, cada uno con un mecanismo distinto.}
\end{frame}

\begin{frame}{Diseño experimental: el RCT}
\begin{block}{4 personas IA, cada una ataca una falla distinta}
\small
\begin{itemize}\setlength{\itemsep}{2pt}
  \item \textbf{T0 --- Control}: Información factual, tono neutral. Corrige directamente las fallas de información
  \item \textbf{T1 --- Friendly Guide}: Calidez conversacional, lenguaje informal. Reduce el costo psicológico de la interacción
  \item \textbf{T2 --- Empathetic Mentor}: Andamiaje psicológico, validación emocional. Ataca la baja autoeficacia
  \item \textbf{T3 --- Inspiring Storyteller}: Narrativas de éxito, modelos a seguir. Eleva los retornos esperados
\end{itemize}
\end{block}

\vspace{0.2cm}
\textbf{Flujo del experimento (3 etapas):}

\medskip
\centering
\begin{tikzpicture}[node distance=2.5cm, auto, >=Stealth,
    box/.style={rectangle, draw=javeriana, thick, fill=javeriana!5, text width=3cm, align=center, minimum height=1cm, font=\small}]
  \node[box] (pre) {1. Encuesta pre\\(baseline, paciencia,\\autoeficacia)};
  \node[box, right of=pre, node distance=4.2cm] (chat) {2. Chat con IA\\(10--15 min,\\persona aleatoria)};
  \node[box, right of=chat, node distance=4.2cm] (post) {3. Encuesta post\\(intención,\\percepciones)};
  \draw[->, thick, javeriana] (pre) -- (chat);
  \draw[->, thick, javeriana] (chat) -- (post);
\end{tikzpicture}

\vspace{0.2cm}
{\small $N = 263$ aleatorizados $\rightarrow$ $n = 148$ completaron (attrición $\approx 44\%$)}
\end{frame}

% ============================================================================
\section{Los datos: Exploración}
% ============================================================================

\begin{frame}[fragile]{Cargar y explorar los datos}
\begin{lstlisting}[basicstyle=\ttfamily\scriptsize]
# Cargar el dataset de-identificado del experimento
datos <- read.csv("datos/rct_student_sample.csv")

# Dimensiones y primeras variables
dim(datos)        # 148 observaciones, 23 variables
names(datos)[1:12]
# "student_id"  "treatment"  "T1"  "T2"  "T3"
# "assigned_group"  "high_patience"  "migrant"
# "parents_college"  "female"  "age"  "gender"

# Muestra por brazo de tratamiento
table(datos$treatment)
# T0_control T1_friendly   T2_mentor T3_narrative
#         37          36          42           33
\end{lstlisting}

\vspace{0.1cm}
\begin{alertblock}{Nota}
\small
Los datos están \textbf{de-identificados}: los IDs son hashes criptográficos (SHA-256), no hay nombres ni emails. Esto es obligatorio en investigación con sujetos humanos.
\end{alertblock}
\end{frame}

\begin{frame}[fragile]{Estadísticas descriptivas}
\begin{lstlisting}[basicstyle=\ttfamily\scriptsize]
# Intención de inscripción (variable principal, escala 0-100)
mean(datos$enrollment_likelihood)  # Baseline: 79.8
sd(datos$enrollment_likelihood)    # SD: 23.0
mean(datos$pt_enrollment_likelihood)  # Post: 82.5

# Intención baseline por brazo de tratamiento
tapply(datos$enrollment_likelihood, datos$treatment, mean)
# T0_control T1_friendly   T2_mentor T3_narrative
#      76.6        78.2        80.0         84.7

# Características de la muestra
mean(datos$female, na.rm = TRUE)          # 51.4% mujeres
mean(datos$migrant, na.rm = TRUE)         # 33.1% migrantes
mean(datos$high_patience, na.rm = TRUE)   # 87.8% pacientes
mean(datos$parents_college, na.rm = TRUE) # 60.8% padres con u.
\end{lstlisting}

\vspace{0.1cm}
{\small \textbf{Observación:} El baseline del grupo T3 (84.7) es más alto que el de T0 (76.6). Diferencia de $\sim$8 puntos \textbf{por azar}. Con solo 148 estudiantes, la aleatorización no garantiza balance perfecto. Por eso necesitamos controles.}
\end{frame}

\begin{frame}[fragile]{Balance: ¿la aleatorización funcionó?}
\begin{exampleblock}{¿Por qué verificar el balance?}
\small
La aleatorización garantiza que, \textit{en expectativa}, los grupos son iguales. Pero con muestras pequeñas puede haber \textbf{desbalance aleatorio}. Verificamos con pruebas F conjuntas.
\end{exampleblock}

\begin{lstlisting}[basicstyle=\ttfamily\scriptsize]
library(fixest)
# Prueba F: cada covariable regresada sobre las dummies
feols(enrollment_likelihood ~ T1 + T2 + T3, data = datos)
# F-test conjunto: F = 0.781, p = 0.507

feols(high_patience ~ T1 + T2 + T3, data = datos)
# F-test conjunto: F = 1.897, p = 0.133
\end{lstlisting}

\vspace{0.1cm}
\textbf{Resultado:} Todas las pruebas F con $p > 0.10$. La aleatorización es válida.

\vspace{0.1cm}
\begin{alertblock}{Diferencias normalizadas (Imbens \& Rubin, 2015)}
\small
Se calculan $d = \frac{\bar{X}_T - \bar{X}_C}{\sqrt{(S_T^2 + S_C^2)/2}}$. El umbral es $|d| < 0.25$. Algunas lo exceden (Storyteller vs Control en baseline: $d = 0.36$), lo que refuerza la necesidad de \textbf{controlar por el baseline}.
\end{alertblock}
\end{frame}

\begin{frame}[fragile]{Visualización: medias crudas por grupo}
\begin{lstlisting}[basicstyle=\ttfamily\scriptsize]
# Medias post-tratamiento por grupo
medias <- tapply(datos$pt_enrollment_likelihood,
                 datos$treatment, mean)
sds    <- tapply(datos$pt_enrollment_likelihood,
                 datos$treatment, sd)
ns     <- tapply(datos$pt_enrollment_likelihood,
                 datos$treatment, length)
se     <- sds / sqrt(ns)

# Grafico de barras con IC 95%
bp <- barplot(medias, col = c("gray70","steelblue","coral","gold"),
              ylim = c(0, 100), ylab = "Intención post (%)",
              main = "Medias crudas por grupo")
arrows(bp, medias - 1.96*se, bp, medias + 1.96*se,
       angle = 90, code = 3, length = 0.05)
\end{lstlisting}

\vspace{0.1cm}
{\small Control: 80.4, Friendly: 80.9, Mentor: 83.0, Storyteller: 85.8. Los intervalos de confianza se solapan ampliamente $\Rightarrow$ las diferencias crudas \textbf{no son significativas}. Pero estas medias no controlan por el baseline.}
\end{frame}

% ============================================================================
\section{El modelo de regresión lineal múltiple}
% ============================================================================

\begin{frame}{El modelo poblacional}
\begin{block}{Modelo}
\[
Y_i = \beta_0 + \beta_1 X_{1i} + \beta_2 X_{2i} + \cdots + \beta_p X_{pi} + \varepsilon_i
\]
donde:
\begin{itemize}
  \item $Y_i$ = variable dependiente (respuesta)
  \item $X_{1i}, X_{2i}, \ldots, X_{pi}$ = variables independientes (predictores, regresores)
  \item $\beta_0$ = intercepto
  \item $\beta_j$ = coeficiente de $X_j$ ($j = 1, 2, \ldots, p$)
  \item $\varepsilon_i$ = error aleatorio
\end{itemize}
\end{block}

\vspace{0.2cm}
\textbf{Interpretación de $\beta_j$:}\\
$\beta_j$ representa el cambio esperado en $Y$ cuando $X_j$ aumenta en una unidad, \textbf{manteniendo constantes todas las demás variables} (ceteris paribus).
\end{frame}

\begin{frame}{Ecuación estimada}
Estimamos los parámetros $\beta_0, \beta_1, \ldots, \beta_p$ usando OLS (mínimos cuadrados ordinarios), obteniendo:
\[
\hat{Y}_i = \hat{\beta}_0 + \hat{\beta}_1 X_{1i} + \hat{\beta}_2 X_{2i} + \cdots + \hat{\beta}_p X_{pi}
\]

\vspace{0.2cm}
\textbf{Residuos:}
\[
e_i = Y_i - \hat{Y}_i
\]

\vspace{0.2cm}
\textbf{OLS minimiza:}
\[
\text{SSE} = \sum_{i=1}^{n} e_i^2 = \sum_{i=1}^{n} (Y_i - \hat{Y}_i)^2
\]

\vspace{0.2cm}
\begin{alertblock}{Nota matricial}
En álgebra matricial: $\boldsymbol{\hat{\beta}} = (\mathbf{X}'\mathbf{X})^{-1}\mathbf{X}'\mathbf{Y}$\\
(No entraremos en detalle, pero R usa esta fórmula internamente)
\end{alertblock}
\end{frame}

\begin{frame}[fragile]{Especificación econométrica del RCT}
\begin{block}{Nuestro modelo (McKenzie, 2012)}
\[
Y_{is} = \alpha + \sum_{k} \beta_k T_{is}^k + \gamma Y_{is}^{\text{pre}} + \mathbf{X}_{is}'\delta + \varepsilon_{is}
\]
\begin{itemize}\setlength{\itemsep}{1pt}
  \item $Y_{is}$ = intención post-tratamiento (0--100) del estudiante $i$
  \item $T_{is}^k$ = dummies de tratamiento (T1, T2, T3; Control es la referencia)
  \item $Y_{is}^{\text{pre}}$ = intención \textbf{baseline} (pre-tratamiento)
  \item $\mathbf{X}_{is}$ = controles: edad, género, migrante, educación padres
  \item $\beta_k$ = \textbf{efecto causal} de cada persona vs.\ Control (ITT)
\end{itemize}
\end{block}

\begin{lstlisting}[basicstyle=\ttfamily\scriptsize]
# En R:
modelo_ols <- lm(pt_enrollment_likelihood ~
                   T1 + T2 + T3 +
                   enrollment_likelihood,
                 data = datos)
\end{lstlisting}
\end{frame}

\begin{frame}[fragile]{Modelo 1: solo dummies de tratamiento}
\begin{lstlisting}[basicstyle=\ttfamily\scriptsize]
modelo_trat <- lm(pt_enrollment_likelihood ~ T1 + T2 + T3,
                  data = datos)
summary(modelo_trat)

# Coefficients:
#             Estimate Std. Error t value Pr(>|t|)
# (Intercept)  80.351      3.519  22.836  < 2e-16 ***
# T1            0.565      5.011   0.113    0.910
# T2            2.649      4.826   0.549    0.584
# T3            5.497      5.125   1.073    0.285
# ---
# Residual standard error: 21.4 on 144 DF
# Multiple R-squared:  0.010, Adjusted R-squared:  -0.011
# F-statistic: 0.467 on 3 and 144 DF,  p-value: 0.706
\end{lstlisting}

\vspace{0.1cm}
{\small \textbf{Resultado:} $R^2 = 0.01$. Ningún tratamiento es significativo. ¿Entonces no sirve? No tan rápido\ldots}
\end{frame}

\begin{frame}[fragile]{Modelo 2: OLS con control baseline (McKenzie, 2012)}
\begin{lstlisting}[basicstyle=\ttfamily\scriptsize]
modelo_ols <- lm(pt_enrollment_likelihood ~ T1 + T2 + T3 +
                   enrollment_likelihood,
                 data = datos)
summary(modelo_ols)

# Coefficients:
#                       Estimate Std. Error t value Pr(>|t|)
# (Intercept)           15.566      2.829    5.50  1.7e-07 ***
# T1                    -0.765      2.060   -0.37    0.711
# T2                    -0.184      1.986   -0.09    0.926
# T3                    -1.306      2.121   -0.62    0.539
# enrollment_likelihood  0.845      0.032   26.64  < 2e-16 ***
# ---
# Residual standard error: 8.80 on 143 DF
# Multiple R-squared:  0.834, Adjusted R-squared:  0.829
# F-statistic: 179.5 on 4 and 143 DF,  p-value: < 2.2e-16
\end{lstlisting}

\vspace{0.1cm}
{\small \textbf{Resultado:} $R^2$ salta de 0.01 a \textbf{0.83}. El baseline explica la mayor parte de la variación. \textbf{Los tratamientos siguen sin ser significativos}, pero ahora tenemos un modelo mucho más preciso. Y los signos cambiaron.}
\end{frame}

\begin{frame}{Interpretación de los coeficientes}
\[
\widehat{Y_{\text{post}}} = 15.57 - 0.77 \times T1 - 0.18 \times T2 - 1.31 \times T3 + 0.85 \times Y_{\text{pre}}
\]

\vspace{0.2cm}
\textbf{Interpretaciones ceteris paribus:}
\begin{itemize}
  \item $\hat{\beta}_{\text{T2}} = -0.18$: Los estudiantes del grupo Mentor reportan 0.18 puntos \textbf{menos} de intención que el control, \textbf{manteniendo constante} la intención baseline. No significativo ($p = 0.93$).

  \item $\hat{\beta}_{\text{baseline}} = 0.85$: Por cada punto adicional de intención previa, la intención post aumenta 0.85 puntos, \textbf{controlando por} el tratamiento asignado. Altamente significativo.
\end{itemize}

\vspace{0.2cm}
\begin{alertblock}{Lección clave}
El control baseline es con frecuencia el mejor predictor del resultado post. Incluirlo es una práctica estándar en economía experimental (McKenzie, \textit{J.\ Dev.\ Econ.}, 2012).
\end{alertblock}
\end{frame}

% ============================================================================
\section{Variables dummy y la trampa}
% ============================================================================

\begin{frame}[fragile]{Variables dummy y la trampa de las dummies}
En nuestro modelo, T1, T2, T3 son \textbf{variables dummy} (0 o 1).

\vspace{0.2cm}
\begin{block}{Categoría de referencia}
T0 (Control) está \textbf{omitido} del modelo. Los coeficientes de T1, T2, T3 se interpretan \textbf{relativos al control}.
\end{block}

\vspace{0.2cm}
\begin{alertblock}{Dummy Variable Trap}
Si incluyéramos T0, T1, T2 y T3 simultáneamente, tendríamos \textbf{multicolinealidad perfecta}: $T0 + T1 + T2 + T3 = 1$ para todo $i$. R elimina automáticamente una, pero debemos saber cuál es la referencia.
\end{alertblock}

\vspace{0.2cm}
\begin{lstlisting}[basicstyle=\ttfamily\scriptsize]
# R lo resuelve automaticamente con factor:
datos$grupo <- factor(datos$assigned_group,
                      labels = c("Control","Friendly","Mentor","Narrative"))
lm(pt_enrollment_likelihood ~ grupo, data = datos)
# Omite grupoControl (la primera categoria)
\end{lstlisting}
\end{frame}

% ============================================================================
\section{$R^2$ y $R^2$ ajustado}
% ============================================================================

\begin{frame}{¿Por qué aprender $R^2$ ajustado?}
\begin{exampleblock}{Porque agregar variables no siempre mejora el modelo}
\small
Si meten 50 variables en la regresión, $R^2$ siempre sube --- pero eso no significa que el modelo sea mejor. Es como agregar filtros a una foto de Instagram: más filtros no la hacen más real. El $R^2$ ajustado penaliza por complejidad innecesaria y les dice si agregar otra variable realmente vale la pena.
\end{exampleblock}
\end{frame}

\begin{frame}{$R^2$ en regresión múltiple}
El $R^2$ se define igual que en regresión simple:
\[
R^2 = \frac{\text{SSR}}{\text{SST}} = 1 - \frac{\text{SSE}}{\text{SST}}
\]

\vspace{0.2cm}
\textbf{Problema:} $R^2$ \textbf{siempre aumenta} (o se mantiene) al agregar variables, incluso si son irrelevantes.

\vspace{0.2cm}
\textbf{Solución:} El $R^2$ ajustado ($R^2_{\text{adj}}$) penaliza por el número de variables:
\[
R^2_{\text{adj}} = 1 - \frac{\text{SSE}/(n - p - 1)}{\text{SST}/(n - 1)} = 1 - \frac{\text{MSE}}{\text{MST}}
\]

\vspace{0.1cm}
\textbf{Propiedades:}
\begin{itemize}
  \item $R^2_{\text{adj}} \leq R^2$
  \item $R^2_{\text{adj}}$ puede disminuir si agregamos una variable irrelevante
  \item Útil para comparar modelos con diferente número de variables
\end{itemize}
\end{frame}

\begin{frame}[fragile]{Comparación: $R^2$ vs $R^2$ ajustado con datos del RCT}
\begin{lstlisting}[basicstyle=\ttfamily\scriptsize]
# Modelo 1: solo tratamientos (3 predictores)
summary(modelo_trat)$r.squared      # 0.010
summary(modelo_trat)$adj.r.squared  # -0.011  (<- negativo!)

# Modelo 2: + baseline (4 predictores)
summary(modelo_ols)$r.squared      # 0.834
summary(modelo_ols)$adj.r.squared  # 0.829

# Modelo 3: + autoeficacia (5 predictores)
modelo_completo <- lm(pt_enrollment_likelihood ~
                        T1 + T2 + T3 +
                        enrollment_likelihood + se_pre,
                      data = datos)
summary(modelo_completo)$r.squared      # 0.835
summary(modelo_completo)$adj.r.squared  # 0.829
\end{lstlisting}

\vspace{0.1cm}
{\small \textbf{Observación:} Agregar \texttt{se\_pre} al Modelo 2 sube $R^2$ de 0.834 a 0.835, pero $R^2_{\text{adj}}$ se mantiene en 0.829. La variable \textbf{no aporta} y el $R^2_{\text{adj}}$ lo detecta.}
\end{frame}

\begin{frame}{¿Cuándo usar $R^2$ ajustado?}
\begin{block}{Regla práctica}
\begin{itemize}
  \item Para \textbf{reportar} el ajuste de un modelo: usar ambos ($R^2$ y $R^2_{\text{adj}}$)
  \item Para \textbf{comparar} modelos con diferente número de variables: usar $R^2_{\text{adj}}$
  \item Para \textbf{seleccionar} el mejor modelo: complementar con AIC/BIC
\end{itemize}
\end{block}

\vspace{0.2cm}
\begin{alertblock}{Precaución}
$R^2_{\text{adj}}$ NO es una probabilidad ni una proporción estricta. Puede ser negativo si el modelo es muy malo --- de hecho, nuestro Modelo 1 (solo tratamientos) tiene $R^2_{\text{adj}} = -0.011$.
\end{alertblock}

\vspace{0.2cm}
\begin{exampleblock}{Interpretación}
Si $R^2_{\text{adj}} = 0.829$, aproximadamente el 83\% de la variabilidad de $Y$ se explica por los predictores, ajustando por el número de variables.
\end{exampleblock}
\end{frame}

% ============================================================================
\section{Pruebas de significancia}
% ============================================================================

\begin{frame}{¿Por qué aprender pruebas de significancia en regresión múltiple?}
\begin{exampleblock}{Porque ahora hay dos preguntas distintas}
\small
\begin{itemize}\setlength{\itemsep}{2pt}
  \item \textbf{Prueba F global}: ¿el modelo \textit{en conjunto} sirve para algo? (Como preguntar: ¿esta playlist tiene al menos una buena canción?)
  \item \textbf{Prueba t individual}: ¿\textit{cada} variable aporta algo, controlando por las demás? (Como preguntar: ¿esta canción específica le aporta a la playlist?)
  \item Pueden tener un modelo útil (F significativo) donde alguna variable individual no sea significativa
\end{itemize}
\end{exampleblock}
\end{frame}

\begin{frame}{Prueba F global}
\begin{block}{Hipótesis}
\begin{itemize}
  \item $H_0$: $\beta_1 = \beta_2 = \cdots = \beta_p = 0$ (ningún predictor es útil)
  \item $H_1$: Al menos un $\beta_j \neq 0$ (al menos un predictor es útil)
\end{itemize}
\end{block}

\vspace{0.2cm}
\textbf{Estadístico:}
\[
F = \frac{\text{MSR}}{\text{MSE}} = \frac{\text{SSR}/p}{\text{SSE}/(n - p - 1)} \sim F_{p, n-p-1} \quad \text{bajo } H_0
\]

\vspace{0.2cm}
\textbf{Regla de decisión:}
Rechazamos $H_0$ si $F > F_{\alpha, p, n-p-1}$ o si p-valor $< \alpha$.

\vspace{0.2cm}
\begin{exampleblock}{Nuestro RCT}
Modelo 2: $F = 179.5$, $p\text{-valor} < 2.2 \times 10^{-16}$. El modelo en conjunto es altamente significativo, aunque los tratamientos individuales no lo son --- \textbf{es el baseline el que empuja}.
\end{exampleblock}
\end{frame}

\begin{frame}{Relación entre $F$ y $R^2$}
El estadístico $F$ se puede expresar en función de $R^2$:
\[
F = \frac{R^2 / p}{(1 - R^2)/(n - p - 1)}
\]

\vspace{0.2cm}
\textbf{Interpretación:}
\begin{itemize}
  \item $R^2$ alto $\Rightarrow$ $F$ alto $\Rightarrow$ modelo significativo
  \item $R^2$ cercano a 0 $\Rightarrow$ $F$ cercano a 0 $\Rightarrow$ modelo no útil
\end{itemize}

\vspace{0.3cm}
\begin{exampleblock}{Ejemplo con nuestros datos}
Modelo 2 con $R^2 = 0.834$, $p = 4$, $n = 148$:
\[
F = \frac{0.834 / 4}{(1 - 0.834) / (148 - 4 - 1)} = \frac{0.2085}{0.00116} \approx 179.5
\]
Con p-valor $< 2.2 \times 10^{-16}$, rechazamos $H_0$.
\end{exampleblock}
\end{frame}

\begin{frame}{Prueba t individual}
\begin{block}{Hipótesis para cada coeficiente}
\begin{itemize}
  \item $H_0$: $\beta_j = 0$ (el predictor $X_j$ NO tiene efecto, controlando por los demás)
  \item $H_1$: $\beta_j \neq 0$ (el predictor $X_j$ sí tiene efecto)
\end{itemize}
\end{block}

\vspace{0.1cm}
\textbf{Estadístico:}
$t_j = \hat{\beta}_j / \text{SE}(\hat{\beta}_j) \sim t_{n - p - 1}$ bajo $H_0$.
Rechazamos $H_0$ si $|t_j| > t_{\alpha/2, n-p-1}$ o si p-valor $< \alpha$.

\vspace{0.1cm}
\begin{alertblock}{Diferencia con regresión simple}
\small En regresión múltiple, la prueba t para $\beta_j$ evalúa el efecto de $X_j$ \textbf{controlando por todas las demás variables}. El resultado puede cambiar drásticamente si agregamos o quitamos predictores.
\end{alertblock}
\end{frame}

\begin{frame}[fragile]{Interpretando el output de \texttt{summary()}}
\begin{lstlisting}[basicstyle=\ttfamily\scriptsize]
summary(modelo_ols)

# Coefficients:
#                       Estimate Std. Error t value Pr(>|t|)
# (Intercept)           15.566      2.829    5.50  1.7e-07 ***
# T1                    -0.765      2.060   -0.37    0.711
# T2                    -0.184      1.986   -0.09    0.926
# T3                    -1.306      2.121   -0.62    0.539
# enrollment_likelihood  0.845      0.032   26.64  < 2e-16 ***
\end{lstlisting}

\vspace{0.2cm}
\textbf{Conclusión:}
\begin{itemize}
  \item Solo \texttt{enrollment\_likelihood} es significativo ($t = 26.64$, $p < 2 \times 10^{-16}$)
  \item Ningún tratamiento es significativo al 5\%
  \item El modelo es significativo \textbf{en conjunto} (F = 179.5) pero por el baseline
\end{itemize}
\end{frame}

\begin{frame}{Significancia conjunta vs individual}
\begin{alertblock}{Discrepancia aparente en nuestro modelo}
\small
\begin{itemize}
  \item Prueba F global: significativa ($p < 2.2 \times 10^{-16}$)
  \item Pruebas t de T1, T2, T3: ninguna significativa
  \item ¿Contradicción? No. La F es significativa porque \texttt{enrollment\_likelihood} es un predictor potente
\end{itemize}
\end{alertblock}

\vspace{0.1cm}
\textbf{Lecciones:}
\begin{itemize}
  \item Un modelo significativo (F) \textbf{no implica} que todos los coeficientes lo sean
  \item La significancia individual depende de qué más está en el modelo
  \item En nuestro caso: la muestra ($n = 148$) puede ser insuficiente para detectar efectos pequeños de tratamiento
\end{itemize}

\vspace{0.1cm}
\begin{block}{Recomendación}
Siempre reportar ambas pruebas (F global e individuales). Si hay discrepancia, investigar multicolinealidad (VIF) y considerar el poder estadístico.
\end{block}
\end{frame}

% ============================================================================
\section{Poder estadístico: ¿tenemos suficientes datos?}
% ============================================================================

\begin{frame}{¿Qué es el poder estadístico?}
\begin{block}{Definición}
El \textbf{poder estadístico} es la probabilidad de detectar un efecto \textit{cuando realmente existe}:
\[
\text{Poder} = P(\text{rechazar } H_0 \mid H_1 \text{ es verdadera}) = 1 - \beta
\]
donde $\beta$ es la probabilidad de \textbf{error Tipo II} (no detectar un efecto real).
\end{block}

\vspace{0.2cm}
\begin{exampleblock}{Analogía: el detector de metales}
\small
El poder es como la sensibilidad de un detector de metales. Si el poder es 80\%, detectarán el tesoro 80 de cada 100 veces que esté ahí. Un poder bajo (ej.\ 30\%) significa que probablemente pasarán por encima del tesoro sin detectarlo.
\end{exampleblock}

\vspace{0.2cm}
\textbf{Convención:} Se busca un poder mínimo del \textbf{80\%} ($\beta \leq 0.20$).

\vspace{0.1cm}
{\small \textbf{¿Por qué importa?} Si nuestros tratamientos no son significativos, ¿es porque realmente no funcionan o porque la muestra es demasiado pequeña para detectarlo?}
\end{frame}

\begin{frame}{Efecto mínimo detectable (MDE)}
\begin{block}{Fórmula}
El MDE es el efecto más pequeño que podemos detectar con poder $1 - \beta$ y significancia $\alpha$:
\[
\text{MDE} = (z_{\alpha/2} + z_{\beta}) \times \sigma \times \sqrt{\frac{1}{n_T} + \frac{1}{n_C}}
\]
Con $\alpha = 0.05$ y poder $= 0.80$: $(z_{0.025} + z_{0.20}) = 1.96 + 0.84 = 2.80$.
\end{block}

\vspace{0.2cm}
\textbf{Sin controles} (usando $\sigma = 21.3$, la SD del outcome):
\[
\text{MDE} = 2.80 \times 21.3 \times \sqrt{\frac{2}{37}} = \textbf{13.9 \text{ pp}}
\]

\textbf{Con control baseline} ($\sigma_{\text{res}} = 21.3 \times \sqrt{1 - 0.83} = 8.8$):
\[
\text{MDE} = 2.80 \times 8.8 \times \sqrt{\frac{2}{37}} = \textbf{5.7 \text{ pp}}
\]
\end{frame}

\begin{frame}{Lección de McKenzie (2012): el poder del baseline}
\begin{alertblock}{El control baseline reduce el MDE 2.4 veces}
¡El MDE pasa de \textbf{13.9 a 5.7 pp}! Incluir el baseline no es decoración: es una herramienta de poder estadístico.
\end{alertblock}

\vspace{0.3cm}
\textbf{¿Por qué funciona?}
\begin{itemize}\setlength{\itemsep}{4pt}
  \item El baseline explica el 83\% de la varianza ($R^2 = 0.83$)
  \item La varianza residual cae de $21.3^2$ a $8.8^2$
  \item Menos ruido $\Rightarrow$ más fácil detectar la señal del tratamiento
  \item Es como usar audífonos con cancelación de ruido para escuchar mejor
\end{itemize}

\vspace{0.3cm}
{\small \textbf{Implicación práctica:} en un RCT, incluir $Y_{\text{pre}}$ puede ser más eficiente que duplicar la muestra.}
\end{frame}

\begin{frame}[fragile]{Poder en nuestro RCT}
\begin{lstlisting}[basicstyle=\ttfamily\scriptsize]
# Calcular poder para distintos tamanios de efecto
power.t.test(n = 37, delta = 5.7, sd = 8.8,
             sig.level = 0.05, type = "two.sample")
# power = 0.803  (justo en el umbral!)

# Nuestros efectos estimados:
# T1 = -0.77, T2 = -0.18, T3 = -1.31
# Todos menores que MDE = 5.7 pp => NO DETECTABLES

# Pero el HTE Mentor x Impaciente = -14.9 pp
power.t.test(n = 18, delta = 14.9, sd = 8.8,
             sig.level = 0.05, type = "two.sample")
# power = 0.99  (detectable incluso con n = 18!)
\end{lstlisting}
\end{frame}

\begin{frame}{Interpretación del análisis de poder}
\begin{block}{¿Qué nos dice el poder sobre nuestros resultados?}
\small
\begin{itemize}\setlength{\itemsep}{4pt}
  \item Efectos promedio de $\sim$1 pp $\Rightarrow$ muy por debajo del MDE (5.7 pp). No podemos saber si son cero o simplemente pequeños
  \item Efecto heterogéneo de $-14.9$ pp $\Rightarrow$ muy por encima del MDE. Detectable incluso con pocos impacienciosos ($n = 18$, poder = 0.99)
\end{itemize}
\end{block}

\vspace{0.3cm}
\begin{alertblock}{Moraleja}
Muestras pequeñas detectan efectos \textbf{grandes} pero no pequeños. Los tratamientos podrían funcionar --- simplemente no tenemos suficiente poder para detectar efectos de 1--2 pp con $n = 37$ por brazo.
\end{alertblock}
\end{frame}

\begin{frame}{Lecciones de poder para su investigación}
\begin{exampleblock}{Reglas prácticas}
\small
\begin{enumerate}\setlength{\itemsep}{4pt}
  \item \textbf{Controles baseline aumentan el poder} --- McKenzie (2012) muestra que en RCTs, incluir $Y_{\text{pre}}$ puede ser más eficiente que duplicar la muestra

  \item \textbf{Muestras pequeñas detectan efectos GRANDES pero no pequeños} --- nuestro RCT con $n = 148$ detecta efectos $\geq 5.7$ pp, no de 1--2 pp

  \item \textbf{Siempre calcular el MDE \textit{antes} de recolectar datos} --- si el MDE es más grande que el efecto esperado, la investigación está condenada

  \item \textbf{Reportar el MDE junto a resultados no significativos} --- ``no encontramos efectos, pero solo podemos descartar efectos mayores a 5.7 pp'' es mucho más informativo que ``no significativo''
\end{enumerate}
\end{exampleblock}

\vspace{0.2cm}
\begin{alertblock}{Error común}
Concluir ``el tratamiento no funciona'' cuando en realidad la muestra no tiene poder para detectarlo. Ausencia de evidencia $\neq$ evidencia de ausencia.
\end{alertblock}
\end{frame}

% ============================================================================
\section{Sesgo por variable omitida (OVB)}
% ============================================================================

\begin{frame}{¿Por qué aprender sobre OVB?}
\begin{exampleblock}{Porque es el error más peligroso en estadística}
\small
\begin{itemize}\setlength{\itemsep}{2pt}
  \item Si su modelo omite una variable importante, los coeficientes están \textbf{sesgados} --- dan respuestas equivocadas
  \item Ejemplo clásico: ``tener iPad mejora las notas'' --- pero en realidad las familias con más recursos compran iPads \textbf{y además} invierten en educación
  \item OVB es la razón por la que no pueden confiar en cualquier correlación que vean en TikTok o Twitter
  \item ¿Y en un experimento? La aleatorización protege del sesgo, pero con muestras pequeñas la protección es imperfecta
\end{itemize}
\end{exampleblock}
\end{frame}

\begin{frame}{¿Qué es el sesgo por variable omitida?}
\begin{block}{Definición}
El \textbf{sesgo por variable omitida} (Omitted Variable Bias, OVB) ocurre cuando:
\begin{enumerate}
  \item Omitimos una variable relevante del modelo
  \item Esa variable está correlacionada con al menos uno de los predictores incluidos
\end{enumerate}
En este caso, los coeficientes estimados están \textbf{sesgados}.
\end{block}

\vspace{0.1cm}
\begin{exampleblock}{Ejemplo}
\small Si omitimos el ingreso familiar del modelo $\text{Notas} = \beta_0 + \beta_1 \times \text{iPad} + \varepsilon$, y el ingreso está correlacionado con tener iPad, entonces $\hat{\beta}_1$ \textbf{sobreestima} el efecto del iPad.
\end{exampleblock}
\end{frame}

\begin{frame}{Fórmula del sesgo}
Si el modelo verdadero es:
\[
Y = \beta_0 + \beta_1 X_1 + \beta_2 X_2 + \varepsilon
\]

Pero estimamos el modelo omitido:
\[
Y = \gamma_0 + \gamma_1 X_1 + u
\]

Entonces el sesgo de $\hat{\gamma}_1$ es:
\[
E(\hat{\gamma}_1) = \beta_1 + \beta_2 \times \frac{\text{Cov}(X_1, X_2)}{\text{Var}(X_1)}
\]

\vspace{0.2cm}
\textbf{Dirección del sesgo:}
\begin{itemize}
  \item Si $\beta_2 > 0$ y $\text{Cov}(X_1, X_2) > 0$ $\Rightarrow$ sesgo positivo (sobreestimación)
  \item Si $\beta_2 > 0$ y $\text{Cov}(X_1, X_2) < 0$ $\Rightarrow$ sesgo negativo (subestimación)
\end{itemize}
\end{frame}

\begin{frame}[fragile]{Ejemplo con el RCT: ¿qué pasa si omitimos el baseline?}
\textbf{Modelo SIN control baseline:}
\begin{lstlisting}[basicstyle=\ttfamily\scriptsize]
coef(modelo_trat)
# (Intercept)        T1        T2        T3
#      80.35      0.57      2.65      5.50

# T3 parece tener efecto +5.50 puntos. Suena prometedor!
\end{lstlisting}

\vspace{0.2cm}
\textbf{Modelo CON control baseline:}
\begin{lstlisting}[basicstyle=\ttfamily\scriptsize]
coef(modelo_ols)
# (Intercept)          T1          T2          T3 enrollment
#      15.57       -0.77       -0.18       -1.31       0.85

# T3 ahora es -1.31. Cambio de signo!
\end{lstlisting}

\vspace{0.2cm}
{\small \textbf{Explicación:} El grupo T3 tenía un baseline más alto \textbf{por azar} (84.7 vs 76.6 del control). Sin controlar, le atribuíamos al tratamiento una ventaja que ya existía antes del experimento.}
\end{frame}

\begin{frame}[fragile]{Verificar el desbalance por azar}
\begin{lstlisting}[basicstyle=\ttfamily\scriptsize]
# Media de enrollment_likelihood (baseline) por grupo
tapply(datos$enrollment_likelihood, datos$treatment, mean)

#  T0_control T1_friendly   T2_mentor T3_narrative
#       76.65       78.22       80.00        84.70

# Con solo 148 estudiantes, hay desbalance aleatorio.
# T3 empezo con ventaja de ~8 puntos sobre el control.
\end{lstlisting}

\vspace{0.2cm}
\begin{alertblock}{Lección para la práctica}
\begin{itemize}
  \item En datos observacionales: OVB crea sesgo real. Hay que controlar por confusores.
  \item En un RCT: la aleatorización protege \textbf{en expectativa}, pero con muestras pequeñas hay desbalance aleatorio.
  \item Controlar por el baseline aumenta la precisión y corrige desbalances (McKenzie, 2012).
\end{itemize}
\end{alertblock}
\end{frame}

\begin{frame}{Estrategias para mitigar OVB}
\begin{block}{En datos observacionales}
\begin{itemize}
  \item Incluir todas las variables confusoras relevantes
  \item Variables instrumentales (IV)
  \item Diseños cuasi-experimentales (DiD, RD)
\end{itemize}
\end{block}

\vspace{0.2cm}
\begin{block}{En datos experimentales (RCT)}
\begin{itemize}
  \item La aleatorización elimina el sesgo \textbf{en promedio}
  \item Incluir controles baseline mejora la precisión (reduce errores estándar)
  \item Esta práctica se llama \textbf{OLS con controles baseline} (McKenzie, 2012)
\end{itemize}
\end{block}

\vspace{0.2cm}
\begin{exampleblock}{Regla de oro}
\small Nunca concluir causalidad solo con regresión observacional. Si ven que ``usar ChatGPT'' se asocia con mejores notas, pregúntense: ¿qué variable omitida podría explicar ambas cosas?
\end{exampleblock}
\end{frame}

% ============================================================================
\section{VIF y multicolinealidad}
% ============================================================================

\begin{frame}{¿Qué es la multicolinealidad?}
\begin{block}{Definición}
La \textbf{multicolinealidad} ocurre cuando dos o más predictores están \textbf{altamente correlacionados} entre sí.
\end{block}

\vspace{0.1cm}
\textbf{Tipos:}
\begin{itemize}
  \item \textbf{Perfecta:} Una variable es combinación lineal exacta de otras (la dummy variable trap)
  \item \textbf{Alta:} Correlación fuerte pero no perfecta (común en datos observacionales)
\end{itemize}

\vspace{0.1cm}
\textbf{Consecuencias:}
\begin{itemize}
  \item Errores estándar \textbf{inflados}, intervalos de confianza muy amplios
  \item Coeficientes inestables, dificultad para detectar significancia individual
\end{itemize}

\vspace{0.2cm}
\begin{alertblock}{Importante}
La multicolinealidad \textbf{no afecta} $R^2$ ni la capacidad de predicción del modelo. Solo afecta la interpretación de los coeficientes individuales.
\end{alertblock}
\end{frame}

\begin{frame}{Factor de Inflación de la Varianza (VIF)}
\begin{block}{Definición}
El \textbf{VIF} (Variance Inflation Factor) mide cuánto se infla la varianza de $\hat{\beta}_j$ debido a la correlación con otros predictores.

Para cada predictor $X_j$:
\[
\text{VIF}_j = \frac{1}{1 - R_j^2}
\]
donde $R_j^2$ es el $R^2$ de la regresión auxiliar de $X_j$ sobre todos los demás predictores.
\end{block}

\vspace{0.2cm}
\textbf{Reglas de pulgar:}
\begin{itemize}
  \item VIF $< 5$ $\Rightarrow$ multicolinealidad \textbf{baja} (aceptable)
  \item $5 \leq$ VIF $< 10$ $\Rightarrow$ multicolinealidad \textbf{moderada} (preocupante)
  \item VIF $\geq 10$ $\Rightarrow$ multicolinealidad \textbf{severa} (requiere acción)
\end{itemize}
\end{frame}

\begin{frame}[fragile]{Calcular VIF en R}
\begin{lstlisting}[basicstyle=\ttfamily\scriptsize]
library(car)

vif(modelo_completo)

#                    T1                    T2
#                  1.51                  1.53
#                    T3 enrollment_likelihood
#                  1.50                  1.05
#                se_pre
#                  1.05

# Todos los VIF < 2 => NO hay multicolinealidad
\end{lstlisting}

\vspace{0.2cm}
\begin{exampleblock}{¿Por qué es tan bajo en un RCT?}
\small
Porque la aleatorización hace que los tratamientos (T1, T2, T3) sean \textbf{casi ortogonales} a las variables baseline. La correlación entre T1 y \texttt{enrollment\_likelihood} es solo $-0.04$. Esto es una ventaja enorme de los experimentos sobre los datos observacionales.
\end{exampleblock}
\end{frame}

\begin{frame}{Soluciones si VIF es alto}
\begin{alertblock}{Opciones}
\begin{enumerate}
  \item \textbf{Eliminar una variable} correlacionada (la menos importante teóricamente)
  \item \textbf{Combinar variables:} crear un índice o promedio (ej: promedio de \texttt{self\_confidence\_1}, \texttt{self\_confidence\_2}, \texttt{self\_confidence\_3} $\rightarrow$ \texttt{se\_pre})
  \item \textbf{Componentes principales (PCA):} reducir dimensionalidad
  \item \textbf{Regularización:} Ridge regression (penaliza coeficientes grandes)
  \item \textbf{Aumentar el tamaño muestral} (reduce errores estándar en general)
\end{enumerate}
\end{alertblock}

\vspace{0.2cm}
\begin{exampleblock}{Nuestro caso}
\small
En el RCT, creamos \texttt{se\_pre} como promedio de \texttt{self\_confidence\_1}, \texttt{self\_confidence\_2}, \texttt{self\_confidence\_3}. Esto es exactamente opción 2: combinar variables correlacionadas en un índice compuesto.
\end{exampleblock}
\end{frame}

% ============================================================================
\section{Interacciones y efectos heterogéneos}
% ============================================================================

\begin{frame}{¿Por qué aprender interacciones?}
\begin{exampleblock}{Porque los efectos no siempre son iguales para todos}
\small
\begin{itemize}\setlength{\itemsep}{2pt}
  \item ¿El efecto del mentor IA es el mismo para un estudiante local que para uno de origen migrante? Probablemente no
  \item ¿El beneficio de un consejero narrativo depende de si la persona ya tenía alta intención de inscribirse? Podría ser
  \item Las interacciones modelan que el efecto de una variable \textbf{depende} de otra --- eso es mucho más realista
\end{itemize}
\end{exampleblock}

\vspace{0.2cm}
\begin{block}{En economía experimental: Efectos Heterogéneos de Tratamiento (HTE)}
\small
¿El efecto promedio nulo ($\hat{\beta} \approx 0$) esconde subgrupos donde el efecto es positivo o negativo? Los HTE son una de las preguntas más importantes en economía aplicada.
\end{block}
\end{frame}

\begin{frame}{¿Qué es una interacción?}
\begin{block}{Definición}
Una \textbf{interacción} ocurre cuando el efecto de una variable $X_1$ sobre $Y$ \textbf{depende del nivel} de otra variable $X_2$.
\end{block}

\vspace{0.2cm}
\textbf{Modelo con interacción:}
\[
Y = \beta_0 + \beta_1 X_1 + \beta_2 X_2 + \beta_3 (X_1 \times X_2) + \varepsilon
\]

\vspace{0.1cm}
\textbf{Interpretación:}
\begin{itemize}
  \item $\beta_1$ = efecto de $X_1$ cuando $X_2 = 0$
  \item $\beta_2$ = efecto de $X_2$ cuando $X_1 = 0$
  \item $\beta_3$ = cómo cambia el efecto de $X_1$ cuando $X_2$ aumenta en 1 unidad
\end{itemize}

\vspace{0.2cm}
El \textbf{efecto marginal de $X_1$} es:
\[
\frac{\partial Y}{\partial X_1} = \beta_1 + \beta_3 X_2
\]
\end{frame}

\begin{frame}[fragile]{Ejemplo: migrant $\times$ tratamiento}
\textbf{Pregunta:} ¿El efecto del mentor IA es diferente para estudiantes de origen migrante?

\begin{lstlisting}[basicstyle=\ttfamily\scriptsize]
modelo_inter <- lm(pt_enrollment_likelihood ~
                     T1 + T2 + T3 +
                     enrollment_likelihood +
                     T1:migrant + T2:migrant + T3:migrant,
                   data = datos)
summary(modelo_inter)

# Coefficients:
#                       Estimate Std. Error t value Pr(>|t|)
# (Intercept)            15.978     2.865    5.58  1.2e-07 ***
# T1                     -0.727     2.345   -0.31    0.757
# T2                      1.368     2.168    0.63    0.529
# T3                     -1.377     2.318   -0.59    0.554
# enrollment_likelihood   0.840     0.032   26.03  < 2e-16 ***
# T1:migrant             -0.081     3.091   -0.03    0.979
# T2:migrant             -5.368     3.016   -1.78    0.077 .
# T3:migrant              0.421     3.436    0.12    0.903
\end{lstlisting}
\end{frame}

\begin{frame}{Interpretación de la interacción T2 $\times$ migrant}
Ecuación para T2 (Mentor):
\[
\hat{Y} = \cdots + 1.37 \times T2 - 5.37 \times (T2 \times \text{migrant}) + \cdots
\]

\vspace{0.2cm}
\textbf{Efecto de T2 según origen:}
\begin{itemize}
  \item No migrante (\texttt{migrant} = 0): efecto $= 1.37$ puntos
  \item Migrante (\texttt{migrant} = 1): efecto $= 1.37 - 5.37 = \mathbf{-4.00}$ puntos
\end{itemize}

\vspace{0.2cm}
\begin{alertblock}{Interpretación}
El Mentor Estructurado parece tener un efecto \textbf{negativo} para estudiantes migrantes ($-4.00$), pero \textbf{positivo} para no migrantes ($+1.37$). La interacción es marginalmente significativa ($p = 0.077$). Con una muestra más grande, podría confirmarse.
\end{alertblock}

\vspace{0.2cm}
{\small \textbf{Nota:} En economía experimental, estos se llaman \textbf{efectos heterogéneos de tratamiento} (Heterogeneous Treatment Effects, HTE).}
\end{frame}

\begin{frame}[fragile]{El hallazgo más importante: Mentor $\times$ Paciencia}
\begin{block}{El efecto promedio nulo esconde heterogeneidad dramática}
\small
¿Recuerdan que 88\% de la muestra es paciente y 12\% impaciente? Veamos qué pasa si interactuamos tratamiento con paciencia:
\end{block}

\begin{lstlisting}[basicstyle=\ttfamily\scriptsize]
modelo_hte <- lm(pt_enrollment_likelihood ~
                   T1 + T2 + T3 + enrollment_likelihood +
                   T1:high_patience + T2:high_patience +
                   T3:high_patience,
                 data = datos)

# Coeficiente clave: T2 (base, impacientes) = -14.9 ***
# Interaccion: T2:high_patience            = +14.0 ***
# Efecto neto para pacientes: -14.9 + 14.0 = -0.9 (n.s.)
\end{lstlisting}

\vspace{0.1cm}
\begin{alertblock}{Resultado central del paper}
\small
\begin{itemize}\setlength{\itemsep}{1pt}
  \item \textbf{Impacientes + Mentor:} $-14.9$ pp ($p < 0.001$). Efecto devastador
  \item \textbf{Pacientes + Mentor:} $\approx -1$ pp (n.s.). Sin efecto
\end{itemize}
\end{alertblock}
\end{frame}

\begin{frame}{Implicaciones del hallazgo HTE}
\begin{block}{La misma intervención, efectos opuestos}
\small
\begin{itemize}\setlength{\itemsep}{4pt}
  \item El Mentor Empático \textbf{daña} a los impacientes ($-14.9$ pp) pero \textbf{no afecta} a los pacientes ($-0.9$ pp, n.s.)
  \item La validación emocional puede reforzar el sesgo al presente en estudiantes impacientes
  \item Un efecto promedio nulo ($\hat{\beta} \approx 0$) escondía esta heterogeneidad dramática
\end{itemize}
\end{block}

\vspace{0.3cm}
\begin{alertblock}{Implicación de política}
Los sistemas ``one-size-fits-all'' de IA pueden dañar a los más vulnerables. Personalizar la intervención según el perfil conductual del estudiante no es un lujo --- es una necesidad.
\end{alertblock}
\end{frame}

\begin{frame}{Tipos de interacción}
\textbf{Continua $\times$ continua:}
\[
Y = \beta_0 + \beta_1 \text{T2} + \beta_2 \text{baseline} + \beta_3 (\text{T2} \times \text{baseline}) + \varepsilon
\]
¿El efecto del Mentor depende del nivel previo de intención?

\vspace{0.3cm}
\textbf{Dummy $\times$ dummy:} (nuestro ejemplo)
\[
Y = \beta_0 + \beta_1 \text{T2} + \beta_2 \text{migrant} + \beta_3 (\text{T2} \times \text{migrant}) + \varepsilon
\]
¿El efecto del Mentor es diferente para migrantes?

\vspace{0.3cm}
\begin{block}{Regla general}
\begin{itemize}
  \item Si $\beta_3 = 0$ $\Rightarrow$ no hay interacción (efecto constante para todos)
  \item Si $\beta_3 \neq 0$ $\Rightarrow$ el efecto de $X_1$ \textbf{varía} según el nivel de $X_2$
\end{itemize}
\end{block}
\end{frame}

% ============================================================================
\section{Robustez: Attrición y Lee bounds}
% ============================================================================

\begin{frame}[fragile]{Attrición: ¿perdimos estudiantes de forma sistemática?}
\begin{block}{El problema}
\small
De 263 estudiantes aleatorizados, solo 148 completaron el experimento (\textbf{attrición del 43.9\%}). Si los que abandonaron lo hicieron \textit{por causa del tratamiento}, nuestros resultados están sesgados.
\end{block}

\begin{lstlisting}[basicstyle=\ttfamily\scriptsize]
# Cargar muestra completa (264 sesiones validas)
# Crear indicador de completitud
# Regresion: completado ~ tratamiento

# Coefficients:
#             Estimate Std. Error t value Pr(>|t|)
# (Intercept)  0.544      0.061    8.94  < 2e-16 ***
# T1           0.056      0.088    0.63    0.527
# T2           0.031      0.084    0.37    0.711
# T3          -0.020      0.088   -0.23    0.817

# Prueba F conjunta: F = 0.285, p = 0.836
\end{lstlisting}

\vspace{0.1cm}
\textbf{Conclusión:} Los tratamientos \textbf{no} causan abandono diferencial. Las tasas de completación son similares: Control 54\%, Friendly 60\%, Mentor 58\%, Storyteller 52\%.
\end{frame}

\begin{frame}{Lee (2009) bounds: ¿qué pasa en el peor caso?}
\begin{block}{La idea}
\small
Incluso sin attrición diferencial, podríamos tener \textbf{attrición selectiva}: quizás los que abandonan son los que habrían tenido peores resultados. Lee (2009) propone ``recortar'' el grupo con más completadores para igualar las tasas, dando cotas sobre el efecto real.
\end{block}

\vspace{0.2cm}
\textbf{Resultados:}
\begin{center}
\small
\begin{tabular}{l c c c c}
\toprule
Brazo & Cota inferior & Cota superior & Efecto crudo & Recorte \\
\midrule
Friendly    & $-1.2$ & $+5.0$ & $+0.6$ & 9.3\% \\
Mentor      & $+1.8$ & $+6.2$ & $+2.6$ & 5.4\% \\
Storyteller & $+3.3$ & $+6.0$ & $+5.5$ & 3.7\% \\
\bottomrule
\end{tabular}
\end{center}

\vspace{0.1cm}
\begin{alertblock}{Hallazgo}
\small
Las cotas del Storyteller \textbf{excluyen el cero} $[3.3, 6.0]$: incluso en el peor escenario de attrición, el Storyteller tiene efecto positivo crudo. Los recortes son modestos (3.7\%--9.3\%), reforzando que la attrición no sesga nuestros resultados principales.
\end{alertblock}
\end{frame}

\begin{frame}{Validez interna: ¿el efecto es causal?}
\begin{block}{Checklist de validez interna}
\begin{enumerate}\setlength{\itemsep}{3pt}
  \item Asignación aleatoria \checkmark{} (verificada con balance de covariables)
  \item No hay attrición diferencial \checkmark{} ($F = 0.285$, $p = 0.836$)
  \item No hay spillovers \checkmark{} (sesiones individuales, no simultáneas)
  \item Lee bounds confirman que attrición no invierte los resultados \checkmark
\end{enumerate}
\end{block}

\vspace{0.3cm}
\textbf{Conclusión:} podemos interpretar los coeficientes como efectos \textbf{causales} (ITT).
\end{frame}

\begin{frame}{Validez externa: ¿se generaliza?}
\begin{alertblock}{Limitaciones honestas}
\small
\begin{itemize}\setlength{\itemsep}{3pt}
  \item Muestra autoseleccionada (participación voluntaria online)
  \item Sesión única de 10--15 minutos (¿qué pasaría con exposición prolongada?)
  \item Intención declarada, no inscripción real
  \item Contexto específico: Flandes, sistema de tracks, jóvenes migrantes
\end{itemize}
\end{alertblock}

\vspace{0.3cm}
{\small En un paper, \textbf{siempre} discutir limitaciones. Esto no debilita el paper --- lo hace más creíble. Un referee que ve limitaciones honestas confía más en los resultados.}
\end{frame}

% ============================================================================
\section{Selección de modelos: AIC y BIC}
% ============================================================================

\begin{frame}{Criterios de información: AIC y BIC}
\begin{block}{Dos criterios populares}
\begin{enumerate}
  \item \textbf{AIC} (Akaike Information Criterion):
  \[
  \text{AIC} = n \ln\!\left(\frac{\text{SSE}}{n}\right) + 2(p + 1)
  \]

  \item \textbf{BIC} (Bayesian Information Criterion):
  \[
  \text{BIC} = n \ln\!\left(\frac{\text{SSE}}{n}\right) + (p + 1) \ln(n)
  \]
\end{enumerate}
\end{block}

\vspace{0.2cm}
{\small \textbf{Interpretación:} Menor AIC/BIC $\Rightarrow$ mejor modelo. BIC penaliza más la complejidad que AIC (cuando $n > 8$). Objetivo predicción: preferir AIC. Objetivo inferencia: preferir BIC.}
\end{frame}

\begin{frame}[fragile]{Calcular AIC y BIC en R}
\begin{lstlisting}[basicstyle=\ttfamily\scriptsize]
AIC(modelo_trat, modelo_ols, modelo_completo, modelo_inter)

#                  df      AIC
# modelo_trat       5  1332.8
# modelo_ols        6  1070.5   <-- MEJOR
# modelo_completo   7  1071.6
# modelo_inter      9  1073.2

BIC(modelo_trat, modelo_ols, modelo_completo, modelo_inter)

#                  df      BIC
# modelo_trat       5  1347.7
# modelo_ols        6  1088.5   <-- MEJOR
# modelo_completo   7  1092.5
# modelo_inter      9  1100.1
\end{lstlisting}

\vspace{0.2cm}
\textbf{Conclusión:} El modelo OLS + baseline (Modelo 2) gana en ambos criterios. Es el más parsimonioso: agregar \texttt{se\_pre} o las interacciones \textbf{no mejora} el modelo.
\end{frame}

\begin{frame}[fragile]{Selección automática con \texttt{step()}}
\begin{lstlisting}[basicstyle=\ttfamily\scriptsize]
# Seleccion automatica hacia atras (backward)
modelo_full <- lm(pt_enrollment_likelihood ~
                    T1 + T2 + T3 +
                    enrollment_likelihood + se_pre + migrant +
                    T1:migrant + T2:migrant + T3:migrant,
                  data = datos)

modelo_step <- step(modelo_full, direction = "backward")

# step() va eliminando variables que no mejoran el AIC:
# 1. Elimina T1:migrant  (p = 0.98)
# 2. Elimina T3:migrant  (p = 0.90)
# 3. Elimina se_pre       (p = 0.35)
# 4. Se detiene: enrollment_likelihood y T2:migrant quedan
\end{lstlisting}

\vspace{0.2cm}
\begin{alertblock}{Precaución}
La selección automática puede llevar a overfitting y p-valores inflados. Úsela como exploración, pero siempre valide con teoría y conocimiento del problema.
\end{alertblock}
\end{frame}

% ============================================================================
\section{Resumen: Anatomía de una investigación}
% ============================================================================

\begin{frame}{Pipeline de una investigación económica}
\begin{block}{Las 10 etapas (lo que hicimos hoy)}
\small
\begin{enumerate}\setlength{\itemsep}{1pt}
  \item \textbf{Pregunta de investigación}: ¿La IA ``cálida'' mejora decisiones educativas?
  \item \textbf{Marco teórico}: Modelo de Capital Humano bajo asimetría de información
  \item \textbf{Diseño experimental}: RCT con 4 brazos, $N = 263$ estudiantes
  \item \textbf{Datos}: Recolección, limpieza, de-identificación, exploración
  \item \textbf{Balance y aleatorización}: Pruebas F, diferencias normalizadas
  \item \textbf{Estimación OLS}: McKenzie (2012), controles baseline, ceteris paribus
  \item \textbf{Significancia y poder}: Pruebas F/t, MDE, interpretar no-significancia
  \item \textbf{Robustez}: OVB, attrición, Lee bounds, validez interna vs.\ externa
  \item \textbf{Heterogeneidad}: Interacciones, HTE por paciencia y migración
  \item \textbf{Conclusiones y limitaciones}: Implicaciones de política + honestidad
\end{enumerate}
\end{block}
{\scriptsize Cada paso tiene una herramienta estadística. Hoy cubrimos los pasos 5--9.}
\end{frame}

\begin{frame}{Problemas comunes y soluciones}
\begin{block}{Sesgo por variable omitida (OVB)}
\begin{itemize}
  \item Omitir variables relevantes correlacionadas con $X$ sesga $\hat{\beta}$
  \item En un RCT, controlar por baseline mejora la precisión (McKenzie, 2012)
\end{itemize}
\end{block}

\begin{block}{Multicolinealidad}
\begin{itemize}
  \item Predictores altamente correlacionados $\Rightarrow$ errores estándar inflados
  \item Detectar con VIF ($> 5$ preocupante, $> 10$ severo)
  \item Solución: eliminar variables redundantes, crear índices, PCA
\end{itemize}
\end{block}

\begin{block}{Interacciones y HTE}
\begin{itemize}
  \item El efecto de $X_1$ depende de $X_2$: $Y \sim X_1 \times X_2$
  \item $\beta_3$ = cómo cambia el efecto de $X_1$ con $X_2$
  \item Un efecto promedio nulo puede esconder heterogeneidad importante
\end{itemize}
\end{block}
\end{frame}

\begin{frame}{Checklist para su examen y futuras investigaciones}
\small
\begin{enumerate}\setlength{\itemsep}{1pt}
  \item \textbf{Exploración inicial:} descriptivas, scatter plots, medias por grupo, balance
  \item \textbf{Ajustar modelo:} \texttt{lm(Y \textasciitilde{} X1 + X2 + ...)}
  \item \textbf{Evaluar ajuste:} $R^2$, $R^2_{\text{adj}}$, prueba F global
  \item \textbf{Significancia individual:} pruebas t, p-valores, IC 95\%
  \item \textbf{Poder estadístico:} MDE, ¿podemos descartar efectos relevantes?
  \item \textbf{Diagnóstico:} Residuos (\texttt{plot}), Multicolinealidad (\texttt{vif})
  \item \textbf{Robustez:} OVB, attrición, Lee bounds
  \item \textbf{Interacciones:} explorar efectos heterogéneos
  \item \textbf{Selección de modelo:} AIC/BIC, $R^2_{\text{adj}}$
  \item \textbf{Reportar:} tabla de coeficientes, MDE, limitaciones, interpretación
\end{enumerate}
\end{frame}

\begin{frame}{Recursos y referencias}
\begin{exampleblock}{Lecturas recomendadas}
\begin{itemize}\setlength{\itemsep}{4pt}
  \item Wooldridge, J. (2019). \textit{Introductory Econometrics}. Capítulos 3--7.
  \item McKenzie, D. (2012). ``Beyond baseline and follow-up.'' \textit{J.\ Dev.\ Econ.}
  \item Lee, D.S. (2009). ``Training, Wages, and Sample Selection.'' \textit{Rev.\ Econ.\ Studies.}
\end{itemize}
\end{exampleblock}

\vspace{0.3cm}
\begin{block}{Datos del experimento}
El dataset de-identificado \texttt{rct\_student\_sample.csv} estará disponible para que practiquen todos los análisis de esta sesión.
\end{block}
\end{frame}

\end{document}
